\section{Core Rules}

The Codex is the physical record of a Runekeeper's Rites. For a traditional Runekeeper, the Codex is a book—vellum pages, leather binding, ink that smells of lampblack and blood. For a Runic Warrior, the Codex is something else entirely.

\subsection{What Is a Codex?}

A Codex is the repository of a Runekeeper's covenant with their Patron. Every Rite the Runekeeper knows is recorded in their Codex. Without a Codex, a Runekeeper cannot invoke Rites—they have no access to the recorded covenants that grant them power.

For a Runic Warrior, the Codex takes a physical form that reflects their Patron and their personal covenant:

\begin{itemize}
\item \textbf{Flesh:} Tattoos, brands, or scarification. The Rites are written on the Runekeeper's skin.
\item \textbf{Stone:} Carved objects—gorgets, shields, staves, standing stones. The Rites are engraved, etched, or carved.
\item \textbf{Steel:} Weapons, armor, or tools. The Rites are engraved on metal. May be sentient.
\end{itemize}

\subsection{Codex Capacity}

A Codex can hold a number of Rites equal to its capacity. Capacity is determined by the Codex's form and quality:

\begin{itemize}
\item \textbf{Minor Codex:} Holds up to 3 Rites. Examples: a small shield, a single tattoo, a simple gorget.
\item \textbf{Standard Codex:} Holds up to 6 Rites. Examples: a full set of flesh-tattoos, a carved staff, an engraved breastplate.
\item \textbf{Major Codex:} Holds up to 9 Rites. Examples: a full body of scarification, a sentient blade, a standing stone.
\end{itemize}

A Runekeeper may only have one active Codex at a time. If they wish to change Codex forms, they must perform a bonding ritual (see below).

\subsection{Creating a Codex}

A Runic Warrior acquires their Codex during character creation or through downtime advancement.

\subsubsection{Initial Codex}

At character creation, a Runic Warrior chooses their Codex form. The Codex is appropriate to their Patron and culture. It has capacity sufficient for their starting Rites.

\subsubsection{Crafting a New Codex}

A Runekeeper may craft a new Codex during downtime. This requires:

\begin{itemize}
\item Appropriate materials (leather, wood, metal, ink, etc.) costing 1 Supply per Rite of capacity.
\item A workspace (forge, studio, quiet room) for the duration.
\item A Craft test (DV 3 for Minor, DV 4 for Standard, DV 5 for Major).
\item Downtime of one day per Rite of capacity.
\end{itemize}

On a success, the Codex is created. On a failure, the materials are ruined. On a critical failure, the Runekeeper marks 1 Fatigue and must acquire new materials.

\subsection{Bonding to a Codex}

A Runekeeper must bond to their Codex to use it. Bonding is a ritual that requires:

\begin{itemize}
\item The Codex itself.
\item A threshold (doorway, crossroads, grave, or altar to their Patron).
\item A witness (the Runekeeper's Thiasos counts).
\item One hour of uninterrupted ritual.
\item Mark 1 Obligation.
\end{itemize}

After bonding, the Runekeeper may invoke Rites from the Codex. If the Codex is lost or destroyed, the Runekeeper loses access to those Rites until they bond to a new Codex.

\subsection{Engraving Rites}

Adding a new Rite to a Codex requires engraving it onto the Codex's surface. For a Runic Warrior, this is a physical act:

\begin{itemize}
\item \textbf{Flesh Codex:} The Rite is tattooed, branded, or scarred onto the skin. Requires a needle, ink (or ash), and a steady hand. The Runekeeper marks 1 Fatigue (the pain of the mark) in addition to the Rite's cost.
\item \textbf{Stone Codex:} The Rite is carved, etched, or engraved. Requires appropriate tools (chisel, hammer, stylus) and a Craft test (DV 3).
\item \textbf{Steel Codex:} The Rite is engraved or etched. Requires a forge or fine tools and a Craft test (DV 4).
\end{itemize}

Engraving a new Rite requires downtime of one day per Tier of the Rite (Low: 1 day, Standard: 2 days, High: 3 days). The Runekeeper must have access to the Rite (from a teacher, a scroll, or Patron instruction).

\subsection{Repairing a Damaged Codex}

A Codex can be damaged. Fire, water, blade, or simple wear.

\begin{itemize}
\item \textbf{Flesh Codex:} Damage heals naturally (rest, Medicine, or magical healing). Damaged tattoos may fade; scars may need to be re-cut. Healing a damaged Flesh Codex requires a Medicine test (DV 3) and one day of rest.
\item \textbf{Stone Codex:} Physical repair requires a Craft test (DV 4) and appropriate materials. A cracked shield can be reforged. A broken staff can be bound with wire.
\item \textbf{Steel Codex:} Requires a Craft test (DV 5) and a forge. A sentient blade may demand a service, a memory, or a promise before it allows repairs.
\end{itemize}

If a Codex is destroyed beyond repair, the Runekeeper loses all Rites engraved on it. They must acquire or craft a new Codex and re-engrave their Rites—a costly and time-consuming process.

\subsection{Losing a Codex}

If a Codex is lost or stolen, the Runekeeper cannot invoke Rites until they recover it or bond to a new Codex.

\begin{itemize}
\item \textbf{Recovery:} The Runekeeper may track the lost Codex (Investigation or Survival test, DV 4-6 depending on circumstances).
\item \textbf{Replacement:} The Runekeeper may craft a new Codex (see above) and bond to it. They lose all Rites engraved on the lost Codex.
\end{itemize}

A sentient Codex (Steel form) may resist being used by anyone other than its bonded Runekeeper. It may refuse to speak, refuse to witness, or actively mislead.

\subsection{The Codex as Witness}

A Codex is not just a record. It is a witness. When a Runic Warrior invokes a Rite, their Codex sees. It remembers. This counts as a witness for Rites that require one (improving Position or reducing DV).

A sentient Codex (Steel form) can also speak. It may offer testimony, confirm an oath, or testify to a covenant. In a hedge-court, a sentient blade's word may carry as much weight as a living witness.

\subsection{Codex and Thiasos}

A sentient Codex may also serve as the Runekeeper's Thiasos (see Chapter: The Thiasos as Ally, Optional Rule: The Sentient Thiasos). This is rare and significant—the Runekeeper's covenant is embodied in a single object that speaks, remembers, and witnesses.

The thread held. The water carried. That is not a Rite. That is grace.
